% =========================================================
% Proof: Preservation of Open and Closed Sets by Bijective Isometries
% Source: notes/isometries/notes-isometries.tex
% Difficulty: Easy
% =========================================================

\subsection*{Preservation of Open and Closed Sets}
\label{prf:isometry-preserves-open}

\begin{remark*}[Return]
$\leftarrow$ \hyperref[prop:isometry-preserves-open]{Back to Proposition in Notes}
\end{remark*}

\begin{proof}
\Claim If $f : X \to Y$ is a bijective isometry then $U \subseteq X$ is
open iff $f(U)$ is open in $Y$.

\Given A bijective isometry $f : X \to Y$ and $U \subseteq X$.

\Goal To prove the biconditional for open sets (the closed case follows by
taking complements and using bijectivity).

\Strategy Use Proposition~\ref{prop:isometry-preserves-balls}:
$f$ maps balls to balls.
So $p \in U$ has a ball $B_r(p) \subseteq U$ iff $f(p) \in f(U)$ has
ball $B_r(f(p)) = f(B_r(p)) \subseteq f(U)$.

% -------------------------------------------------------
% YOUR PROOF HERE
% -------------------------------------------------------

\end{proof}

\begin{remark*}[Proof shape]
Ball-preservation translates the openness condition directly.
\end{remark*}

\begin{remark*}[Dependencies]
\begin{itemize}
  \item Proposition~\ref{prop:isometry-preserves-balls}: $f(B_r(x)) = B_r(f(x))$.
  \item Definition~\ref{def:open-set}: open via balls.
  \item Bijectivity for the closed case: $f(X \setminus U) = Y \setminus f(U)$.
\end{itemize}
\end{remark*}
